\section{Introducción}

\noindent En este informe se estudia el Multi-Source Capacitated Facility Location Problem with Customer Incompatibilities (MS-CFLP-CI), una variante del problema de localización de instalaciones capacitadas aplicada a redes logísticas. El problema consiste en decidir qué bodegas abrir y cómo distribuir la demanda de los clientes entre ellas, considerando costos de apertura, costos de operación, capacidades limitadas e incompatibilidades entre ciertos pares de clientes \cite{klose2005facility,maia2023metaheuristic}.\\

\noindent La motivación de este estudio se relaciona con escenarios logísticos donde no basta con minimizar costos, sino que también se deben cumplir restricciones operacionales. En particular, ciertos clientes o productos pueden no ser compatibles dentro de una misma bodega debido a condiciones de almacenamiento, seguridad, regulación o manejo operacional. Por ello, una solución factible debe satisfacer completamente la demanda de cada cliente, respetar la capacidad disponible en cada bodega y evitar que clientes incompatibles sean atendidos por una misma instalación.\\

\noindent En la variante multi-source, un cliente puede recibir productos desde más de una bodega, lo que entrega mayor flexibilidad para utilizar las capacidades disponibles. Sin embargo, esta característica también aumenta la complejidad del problema, ya que las incompatibilidades deben verificarse en cada bodega donde exista una asignación positiva de demanda \cite{mansini2022kernel,ceschia2024multineighborhood}.\\

\noindent Para abordar el problema, se presenta una representación computacional de las soluciones, junto con una estrategia heurística basada en Hill Climbing + Restart. Este acercamiento permite explorar soluciones vecinas, mejorar progresivamente el costo total y reducir la probabilidad de quedar atrapado en óptimos locales mediante reinicios. Además, se consideran casos de prueba de distinto tamaño para analizar el comportamiento del algoritmo en términos de factibilidad, calidad de solución y desempeño computacional.\\

\noindent El propósito del documento es entregar una base conceptual, formal y experimental para el estudio del MS-CFLP-CI. Para ello, se define el problema y sus componentes principales, se revisa el estado del arte asociado a enfoques exactos, heurísticos, metaheurísticos y matheurísticos, se presenta un modelo matemático, se describe la representación utilizada y se expone el procedimiento de búsqueda local implementado. Finalmente, se analizan los resultados obtenidos en los casos de prueba y se presentan las conclusiones del trabajo.